%============================ TESTS ET VALIDATION ===========================
\chapter{Tests et validation}

\section{Stratégie de validation}

La validation du prototype a été réalisée de manière \textbf{manuelle et
systématique}, module par module, puis de bout en bout. En l'absence de tests
automatisés à ce stade (voir limites), l'accent a été mis sur des vérifications
reproductibles : build, contrôle de santé de l'API, exécution des endpoints via
Swagger, et parcours utilisateur complet dans le navigateur.

\section{Vérifications réalisées}

Le tableau~\ref{tab:tests} récapitule les principales vérifications et leur
résultat.

\begin{table}[H]
\centering
\caption{Synthèse des tests et validations}
\label{tab:tests}
\small
\begin{tabularx}{\textwidth}{@{}p{5.2cm}L p{2cm}@{}}
\toprule
\textbf{Vérification} & \textbf{Détail} & \textbf{Résultat} \\
\midrule
Build du frontend & \texttt{npm run build} sans erreur & Validé \\
Analyse statique & \texttt{npm run lint} sans avertissement & Validé \\
Santé du backend & \texttt{GET /api/health} renvoie le statut \texttt{ok} & Validé \\
API Swagger & Endpoints listés et exécutables depuis \texttt{/docs} & Validé \\
Navigation & Menu adapté au rôle, fil d'Ariane correct & Validé \\
Connexion / déconnexion & Ouverture et fermeture de session & Validé \\
Routes protégées & Redirection si non authentifié ou rôle non autorisé & Validé \\
Création de campagne & Parcours complet et enregistrement du brouillon & Validé \\
Import des magasins & Validation des lignes et détection d'erreurs & Validé \\
Carte de geofencing & Marqueurs et rayons affichés correctement & Validé \\
Génération DCO & Variantes par magasin et \emph{landing pages} simulées & Validé \\
Filtres de reporting & Indicateurs et graphiques recalculés & Validé \\
Export CSV & Fichier lisible dans Excel (accents, colonnes) & Validé \\
Paramètres du compte & Mise à jour via \texttt{PATCH /api/account-settings} & Validé \\
Validation fonctionnelle finale & Parcours de bout en bout cohérent & Validé \\
\bottomrule
\end{tabularx}
\end{table}

\section{Contrôle de non-régression}

À chaque nouvelle étape, les modules précédemment validés ont été parcourus
de nouveau afin de détecter d'éventuelles régressions. Cette vigilance a
notamment permis de corriger, lors de l'intégration de bout en bout, des
détails de navigation (par exemple des liens externes de démonstration rendus
volontairement non cliquables pour éviter d'ouvrir des pages inexistantes).

\section{Cas d'erreur couverts}

Au-delà des cas nominaux, plusieurs cas d'erreur ont été vérifiés côté API :
identifiant inconnu (404), e-mail déjà utilisé lors de la création d'un
utilisateur (409), champ invalide ou manquant (422), fichier d'import au
format non supporté ou illisible (400). Ces réponses sont explicites et
exploitées par l'interface pour afficher des messages clairs.

\section{Limite de la démarche}

La validation reste \textbf{manuelle} : elle ne remplace pas une suite de
\textbf{tests automatisés} (unitaires et de bout en bout), dont la mise en
place figure parmi les perspectives du projet. La procédure de validation est
toutefois documentée et reproductible (voir le chapitre suivant).
